

Now, under the assumption of realizable case, 
we explain how one can estimate the parameters $p$ and $\sigma$
given tests results from a homogeneous population.
Surprisingly, we discover that parameter recovery
in this model does not require any ground truth labels
indicating whether an employee is skilled or unskilled.
% 
%
We use Hoeffding’s inequality to bound the absolute difference 
between the estimated parameters and the true parameters 
by choosing $\delta$ as the wanted upper bound 
and solving for the number of samples or $\epsilon$.
\begin{lemma}[Hoeffding’s inequality]%http://www.stat.cmu.edu/~larry/=stat705/Lecture2.pdf
Let $y_1,\dots,y_m$ be $\sigma^2-$sub--gaussian random variables. 
Then, for any $\epsilon>0$,
\[
\Pr\left[\left|\frac{1}{m}\sum_{i=1}^m y_i-\E[y_i]\right|\geq \epsilon\right]\leq 2e^{-m\epsilon^2/2\sigma^2}.
\]
If $y_1,\dots,y_m$ are Bernoulli random variables with parameter $p$,
\[
\Pr\left[\left|\frac{1}{m}\sum_{i=1}^m y_i-p\right|\geq \epsilon\right]\leq 2e^{-2m\epsilon^2}.
\]
\end{lemma}
%
We start by estimating $\sigma$ and then use it to derive an estimate for $p$. 
The estimated parameters are denoted by $\hat{\sigma}$ and $\hat{p}$.
Notice that in order to have any information regarding the true value of $\sigma$, 
we need to have candidates with at least two tests. 
Hence, from now on we assume exactly that, i.e., 
$\forall_i \pi_{\text{Greedy}}(\hat{y}_{i,1})=more$ for dynamic policies 
and $\tau\geq2$ for fixed number of tests policies. 

Now, in both policies we have showed that the optimal rule 
is to reject candidates that fail their first test. 
Therefore inconsistencies between the first two tests 
are seen only in cases where $\hat{y}_{i,1}=1,\hat{y}_{i,2}=0$.

Let $c$ be the number of inconsistencies in the first two tests, 
i.e., $c=|\{(\hat{y}_{i,1},\hat{y}_{i,2}):y_{i,1}\ne y_{i,2}\}|$, 
and let $m$ be the number of candidates with at least two tests. 
Since $c$ is generated by sampling $m$ times,
the distribution $Br((\frac{1+\sigma}{2})(\frac{1-\sigma}{2}))=Br(\frac{1-\sigma^2}{4})$ and we can estimate $\sigma$ as stated in the next theorem:
\begin{theorem}
If we have results from $m\geq \frac{1}{2\epsilon^2} \ln \frac{2}{\delta}$ candidates, 
by using $\hat{\sigma}=\sqrt{1-4\frac{c}{m}}$,
then with probability $1-\delta$ we have that $|\hat{\sigma}-\sigma|\leq \epsilon$.
\end{theorem}
Having an estimation of the parameter $\hat{\sigma}$, 
we can calculate the estimated $p$ as follows:
%
Let $p_{\hat{y}_{*,1}=1}:=\frac{\sum_i\mathbb{I}(\hat{y}_{i,1}=1)}{m}$ be 
the percentage of positive first tests. 
%
Since this number is generated by the distribution $Br(\frac{1}{2}(p(1+\sigma)+(1-p)(1-\sigma)))=Br(\frac{1}{2}+(2p-1)\frac{\sigma}{2})$, 
we can estimate $\hat{p}$ using the estimated value of $\hat{\sigma}$.

\begin{theorem}
If we have results from $m\geq \frac{1}{2\epsilon^2} \ln \frac{2}{\delta}$ candidates, 
by using $\hat{p}=\frac{2(p_{y_{*,1}=1}-1)+\hat{\sigma}}{\hat{{\sigma}}}$,
we get that with probability $1-\delta$ we have that $|\hat{p}-p|\leq 2\epsilon$.
\end{theorem}
Under the Gaussian screening model,
the parameter estimation is also straightforward 
(assuming realizability) without access to the true skill level of the employees. 
We start by looking at a single candidate, $i$. 
Each of his test results, $\hat{y}_{i,j}$ is generated 
from a conditional distribution $P(Y_{i}|Q_i=q_i)$ 
which is a Gaussian with mean $q_i$ and variance $\sigma_\eta^2$. 
Since this variance is common among all the candidates, 
we can simply average the estimated variance of every candidate 
to get an approximation for $\sigma_\eta^2$.
Suppose $\hat{y}_{i,1},\dots,\hat{y}_{i,n}$ is a sequence of $n$ i.i.d tests of candidate $i$, 
and let $\boldsymbol{y_i}=\frac{1}{n}\sum_{j=1}^n y_{i,j}$ 
be the empirical mean of candidate $i$'s tests.

The following theorem is a result from Hoeffding’s Inequality, 
in which we use to bound the error of our estimated parameters.
\begin{theorem}
By using the following as estimators for Gaussian parameters $\hat{\mu}_Q=\frac{1}{m}\sum_{i=1}^m\boldsymbol{y_i}$,
$\hat{\sigma}_{\eta}^2=\frac{1}{m}\sum_{i=1}^m\frac{1}{n}\sum_{j=1}^{n} (y_{i,j}-\boldsymbol{y_i})^2$ and $\hat{\sigma}_{Q}^2=\frac{1}{m}\sum_{i=1}^m (\hat{\mu}_Q-\boldsymbol{y_i})^2$
(notice that $\E[\hat{\sigma}_{\eta}^2]={\sigma}_{\eta}^2$ and  $\E[\hat{\sigma}_{Q}^2]={\sigma}_{Q}^2$), the difference between each parameter and it's estimator is bounded by $O(\sqrt{\frac{1}{m}\ln(\frac{1}{\delta})})$.
\end{theorem}